\documentclass[presentation]{beamer}
\usepackage{graphicx}
\usepackage{hyperref}

\title{Performance and Calibration of the ATLAS Tile Calorimeter}
\author{Tade\'a\v{s} Petr\'u\\ Charles University\\ On behalf of the ATLAS Collaboration}
\date{26 August - 4 September 2024, Kolymbari\\ XIII International Conference on New Frontiers in Physics}

\begin{document}

\begin{frame}
    \titlepage
\end{frame}

\begin{frame}{Introduction}
    \begin{itemize}
        \item The Tile Calorimeter (TileCal) is the hadronic calorimeter of the ATLAS detector.
        \item Covers central pseudorapidity range of $|\eta| < 1.7$.
        \item Identifies and measures energy of hadrons, jets, tau leptons.
        \item Provides input for L1Calo Trigger and muon identification.
    \end{itemize}
\end{frame}

\begin{frame}{Signal Reconstruction}
    \begin{itemize}
        \item TileCal consists of plastic scintillator tiles and steel absorbers.
        \item Light from scintillating tiles is collected by WLS fibres to PMTs.
        \item Signal shaped and digitized in front-end electronics.
        \item Two readout gains (high and low, 64:1 ratio).
        \item Signal sampled every 25 ns, reconstructed using optimal filtering algorithm.
    \end{itemize}
\end{frame}

\begin{frame}{TileCal Readout Geometry}
    \begin{itemize}
        \item WLS fibres grouped to PMTs to form readout cell geometry.
        \item Two PMTs per readout cell, total of 9852 channels and 5182 cells.
        \item TileCal has four partitions: LBA, LBC, EBA, EBC.
        \item Readout cells divided into three layers in LB and EB.
        \item Each barrel has 64 modules in azimuthal direction.
    \end{itemize}
\end{frame}

\begin{frame}{Energy Calibration}
    \begin{itemize}
        \item Several calibration systems:
        \begin{itemize}
            \item Cesium system
            \item Laser calibration system
            \item Charge injection system
            \item Minimum-bias system
        \end{itemize}
        \item Energy in GeV calculated as:
        \begin{equation}
            E [GeV] = \frac{A[ADC]}{C_{ADC \to pC} \times C_{pC \to GeV} \times C_{Cs} \times C_{MB} \times C_{Las}}
        \end{equation}
    \end{itemize}
\end{frame}

\begin{frame}{Charge Injection System}
    \begin{itemize}
        \item CIS calibrates front-end electronics by injecting defined charge.
        \item Covers full dynamic range.
        \item Estimates conversion constant from ADC to pC with precision of $\sim0.7\%$.
    \end{itemize}
\end{frame}

\begin{frame}{Cesium System}
    \begin{itemize}
        \item Calibrates entire readout chain (optical components, PMTs).
        \item $^{137}$Cs capsule ($\gamma$-source, $E_\gamma = 0.662$ MeV) measures tile responses.
        \item Cesium scans performed monthly with accuracy of $\sim0.3\%$.
    \end{itemize}
\end{frame}

\begin{frame}{Laser System}
    \begin{itemize}
        \item Produces 10 ns light pulses at 532 nm distributed to all PMTs.
        \item Calibrates readout variations and PMT gain changes.
        \item Runs every 2-3 days, precision at $\sim0.5\%$.
    \end{itemize}
\end{frame}

\begin{frame}{Minimum Bias System}
    \begin{itemize}
        \item Measures PMT current proportional to LHC luminosity.
        \item Data collected with integrator readout system.
        \item Calibration constant $C_{MB}$ used for energy reconstruction.
    \end{itemize}
\end{frame}

\begin{frame}{Time Calibration and Performance}
    \begin{itemize}
        \item Precise timing calibration critical for optimal filtering algorithm.
        \item Time-of-flight measurements used for long-lived particle searches.
        \item Calibration corrects timing jumps (threshold reduced to 1 ns).
    \end{itemize}
\end{frame}

\begin{frame}{Performance - Isolated Muons}
    \begin{itemize}
        \item Studied using isolated muons from $W \to \mu \nu$ decays.
        \item Response quantified using $\Delta E / \Delta x$ distributions.
        \item Double ratio given by:
        \begin{equation}
            R = \frac{\langle\Delta E / \Delta x\rangle_{\text{Data}}}{\langle\Delta E / \Delta x\rangle_{\text{MC}}}
        \end{equation}
    \end{itemize}
\end{frame}

\begin{frame}{Performance - Single Isolated Hadrons}
    \begin{itemize}
        \item Isolated hadrons used to measure deposited energy.
        \item Particles selected with $p < 20$ GeV.
        \item Response characterized by $R = E/p$.
        \item Run 2 results: $\langle E/p \rangle = 0.5896 \pm 0.0001$ (data), $0.593 \pm 0.001$ (MC).
    \end{itemize}
\end{frame}

\begin{frame}{Summary}
    \begin{itemize}
        \item TileCal continues operations in Run 3.
        \item Calibration ensures good performance.
        \item Performance studies validate procedures using isolated muons, hadrons, and jets.
        \item First Run 3 results presented, more to come.
    \end{itemize}
\end{frame}

\begin{frame}{References}
    \begin{itemize}
        \item \href{https://twiki.cern.ch/twiki/bin/view/AtlasPublic/ApprovedPlotsTile}{Public plots for TileCal}
        \item \href{https://twiki.cern.ch/twiki/bin/view/AtlasPublic/TileCaloPublicResults}{Collision Data Public Plots}
        \item \href{https://arxiv.org/abs/2401.16034}{ATLAS TileCal Performance in LHC Run 2}
    \end{itemize}
\end{frame}

\end{document}
